\documentclass[11pt,a4paper]{article}
\usepackage[utf8]{inputenc}
\usepackage[spanish,es-noquot]{babel}
\usepackage{amsmath,amssymb}
\usepackage{geometry}
\usepackage{xcolor}
\usepackage{enumitem}
\usepackage{tikz}
\usetikzlibrary{arrows.meta}
\usetikzlibrary{patterns,arrows.meta}

% Desactivar atajos conflictivos de babel en el preámbulo
\spanishdeactivate{<>}

% Configuración de márgenes
\geometry{top=2cm, bottom=2cm, left=2cm, right=2cm}

% Definición del color usado en el documento
\definecolor{numcolor}{rgb}{0.1, 0.4, 0.7}
\definecolor{solcolor}{rgb}{0.0, 0.5, 0.2}

\begin{document}

% Asegura la desactivación de los atajos al iniciar el cuerpo del documento
\shorthandoff{<>}

\begin{center}
    {\LARGE \textbf{Solucionario: Ficha de Repaso}}\\[0.2cm]
    {\large \textbf{Polinomios y Fracciones Algebraicas}}\\[0.4cm]
\end{center}


\begin{enumerate}[label=\textbf{\arabic*.}, leftmargin=*]

    \item Dados $P(x) = x - 3$, $Q(x) = x - 4$ y $R(x) = x + 5$, realiza las siguientes operaciones:
        \begin{enumerate}[label=\textit{\alph*})]
            \item Halla $M(x) = P(x) \cdot Q(x) \cdot R(x)$.
            
            \noindent{\color{solcolor}\textbf{Solución:}} Multiplicamos los polinomios paso a paso:
            $$M(x) = (x - 3)(x - 4)(x + 5) = (x^2 - 7x + 12)(x + 5)$$
            $$M(x) = x^3 + 5x^2 - 7x^2 - 35x + 12x + 60 = \mathbf{x^3 - 2x^2 - 23x + 60}$$

            \item Descompón en factores $M(x)$.
            
            \noindent{\color{solcolor}\textbf{Solución:}} Como $M(x)$ proviene directamente del producto de factores dados al inicio, su descomposición factorial es inmediata:
            $$\mathbf{M(x) = (x - 3)(x - 4)(x + 5)}$$

            \item ¿Cuáles son las raíces de $M(x)$?
            
            \noindent{\color{solcolor}\textbf{Solución:}} Las raíces son los valores que anulan cada factor: 
            $$\mathbf{x = 3}, \quad \mathbf{x = 4} \quad \text{y} \quad \mathbf{x = -5}$$

            \item ¿Para qué valores de $x$ se anula el polinomio $M(x)$?
            
            \noindent{\color{solcolor}\textbf{Solución:}} El polinomio se anula exactamente en sus raíces, es decir, para $\mathbf{x \in \{-5, 3, 4\}}$.
        \end{enumerate}

    \item \textbf{Halla el cociente $C(x)$ y el resto $R(x)$ de las siguientes divisiones de polinomios:}
        \begin{enumerate}[label=\alph*), itemsep=1em]
            \item $(-3x^4 - 5x^3 + 4x^2 - 1) : (x^2 + 2)$
            
            \noindent{\color{solcolor}\textbf{Solución:}} Realizando la división polinómica:
            $$\mathbf{C(x) = -3x^2 - 5x + 10}, \quad \mathbf{R(x) = 10x - 21}$$

            \item $(2x^5 - 4x^3 + 6x^2 - 7x + 2) : (x^2 - 3x + 1)$
            
            \noindent{\color{solcolor}\textbf{Solución:}} 
            $$\mathbf{C(x) = 2x^3 + 6x^2 + 8x + 16}, \quad \mathbf{R(x) = 17x - 14}$$

            \item $(4x^4 - 6x^3 + 2x - 5) : (2x^2 + 1)$
            
            \noindent{\color{solcolor}\textbf{Solución:}} 
            $$\mathbf{C(x) = 2x^2 - 3x - 1}, \quad \mathbf{R(x) = 5x - 4}$$
        \end{enumerate}

    \item \textbf{Mediante la regla de Ruffini, calcula el cociente y el resto de las siguientes divisiones:}
        \begin{enumerate}[label=\textit{\alph*})]
            \item $(x^2 - 7x + 6) : (x - 1) =$
            
            \noindent{\color{solcolor}\textbf{Solución:}} Aplicando Ruffini con raíz $x = 1$:
            \begin{center}
            \begin{tabular}{r|rrr}
              1 & 1 & -7 & 6 \\
              & & 1 & -6 \\
              \hline
              & 1 & -6 & \multicolumn{1}{|r}{0}
            \end{tabular}
            \end{center}
            $$\mathbf{C(x) = x - 6}, \quad \mathbf{R = 0}$$

            \item $(2x^3 + 3x^2 - x - 5) : (x + 2) =$
            
            \noindent{\color{solcolor}\textbf{Solución:}} Aplicando Ruffini con raíz $x = -2$:
            \begin{center}
            \begin{tabular}{r|rrrr}
              -2 & 2 & 3 & -1 & -5 \\
              & & -4 & 2 & -2 \\
              \hline
              & 2 & -1 & 1 & \multicolumn{1}{|r}{-7}
            \end{tabular}
            \end{center}
            $$\mathbf{C(x) = 2x^2 - x + 1}, \quad \mathbf{R = -7}$$

            \item $(x^4 - 8) : (x - 2) =$
            
            \noindent{\color{solcolor}\textbf{Solución:}} Completando los coeficientes intermedios con ceros ($x^4 + 0x^3 + 0x^2 + 0x - 8$) y raíz $x = 2$:
            \begin{center}
            \begin{tabular}{r|rrrrr}
              2 & 1 & 0 & 0 & 0 & -8 \\
              & & 2 & 4 & 8 & 16 \\
              \hline
              & 1 & 2 & 4 & 8 & \multicolumn{1}{|r}{8}
            \end{tabular}
            \end{center}
            $$\mathbf{C(x) = x^3 + 2x^2 + 4x + 8}, \quad \mathbf{R = 8}$$
        \end{enumerate}

    \item \textbf{Descompón en factores los siguientes polinomios:}
        \begin{enumerate}[label=\alph*), itemsep=0.8em]
            \item $x^4 - 12x^3 + 36x^2 = x^2(x^2 - 12x + 36) = \mathbf{x^2(x - 6)^2}$
            \item $2x^4 + 5x^3 - 4x^2 - 3x = x(2x^3 + 5x^2 - 4x - 3)$. Por Ruffini con raíces $1$ y $-1$:
            $\mathbf{x(x - 1)(x + 1)(2x + 3)}$
            \item $x^4 - 16 = (x^2 - 4)(x^2 + 4) = \mathbf{(x - 2)(x + 2)(x^2 + 4)}$
            \item $3x^3 - 12x^2 + 9x = 3x(x^2 - 4x + 3) = \mathbf{3x(x - 1)(x - 3)}$
            \item $x^4 - 2x^3 - 7x^2 + 8x + 12$. Probando divisores por Ruffini (raíces $-1, 2, 3, -2$):
            $\mathbf{(x + 1)(x - 2)(x - 3)(x + 2)}$
            \item $4x - 24x^2 + 16x^3 = 4x(4x^2 - 6x + 1)$ (o sacando factor común): $\mathbf{4x(4x^2 - 6x + 1)}$
            \item $3(x - 1) + 8(x - 1)^2 - 5(x - 1)(x + 1) = (x - 1)[3 + 8(x - 1) - 5(x + 1)] = \mathbf{(x - 1)(3x - 6)} = \mathbf{3(x - 1)(x - 2)}$
            \item $(x + 3)^2 - 4(x + 3)(x - 3) + 6(x + 3) = (x + 3)[(x + 3) - 4(x - 3) + 6] = \mathbf{(x + 3)(-3x + 21)}$
            \item $x^2 + 4x + 4 = \mathbf{(x + 2)^2}$
            \item $x^2 - 9 = \mathbf{(x - 3)(x + 3)}$
            \item $25x^2 - 125 = 25(x^2 - 5) = \mathbf{25(x - \sqrt{5})(x + \sqrt{5})}$
        \end{enumerate}

    \item \textbf{Calcula el valor del parámetro {\color{numcolor}$m$} en cada uno de los siguientes casos:}
        \begin{enumerate}[label=\alph*), itemsep=0.8em]
            \item Para que el polinomio $P(x) = 7x^3 - mx^2 + 3x - 2$ sea divisible por $x + 1$.\\
            \noindent{\color{solcolor}\textbf{Solución:}} Por el teorema del resto, $P(-1) = 0$:
            $$7(-1)^3 - m(-1)^2 + 3(-1) - 2 = 0 \implies -7 - m - 3 - 2 = 0 \implies \mathbf{m = -12}$$
            
            \item Para que la división $(2x^3 + mx^2 - 5x + 6) : (x - 2)$ tenga un resto igual a $12$.\\
            \noindent{\color{solcolor}\textbf{Solución:}} Evaluamos en $x = 2$ e igualamos a $12$:
            $$2(2)^3 + m(2)^2 - 5(2) + 6 = 12 \implies 16 + 4m - 10 + 6 = 12 \implies 4m + 12 = 12 \implies \mathbf{m = 0}$$
            
            \item Para que $x = 3$ sea una raíz del polinomio $Q(x) = x^3 - 4x^2 + mx - 9$.\\
            \noindent{\color{solcolor}\textbf{Solución:}} Evaluamos $Q(3) = 0$:
            $$3^3 - 4(3)^2 + m(3) - 9 = 0 \implies 27 - 36 + 3m - 9 = 0 \implies 3m - 18 = 0 \implies \mathbf{m = 6}$$
        \end{enumerate}

    \item \textbf{Simplifica las siguientes fracciones algebraicas:}
        \begin{enumerate}[label=\alph*), itemsep=1em]
            \item $\dfrac{x^3 - 4x}{x^3 + x^2 - 2x} = \dfrac{x(x-2)(x+2)}{x(x-1)(x+2)} = \mathbf{\dfrac{x - 2}{x - 1}}$
            \item $\dfrac{x^2 + 2x - 3}{x^3 + 6x^2 + 5x - 12} = \dfrac{(x-1)(x+3)}{(x-1)(x+3)(x+4)} = \mathbf{\dfrac{1}{x + 4}}$
            \item $\dfrac{2x^3 - 8x}{x^3 - 2x^2 - 4x + 8} = \dfrac{2x(x-2)(x+2)}{(x-2)^2(x+2)} = \mathbf{\dfrac{2x}{x - 2}}$
            \item $\dfrac{x^2 - 9}{x^2 - 6x + 9} = \dfrac{(x-3)(x+3)}{(x-3)^2} = \mathbf{\dfrac{x + 3}{x - 3}}$
        \end{enumerate}

    \item \textbf{Efectúa las operaciones y simplifica el resultado al máximo:}
        \begin{enumerate}[label=\alph*), itemsep=1.2em]
            \item $\dfrac{2x^2}{x - 3} : \dfrac{8}{x^3 - 3x^2} = \dfrac{2x^2 \cdot x^2(x - 3)}{8(x - 3)} = \dfrac{2x^4(x - 3)}{8(x - 3)} = \mathbf{\dfrac{x^4}{4}}$
            \item $\dfrac{x^2 - 6}{(x - 2)^2} - \dfrac{x - 3}{x - 2} = \dfrac{(x^2 - 6) - (x - 3)(x - 2)}{(x - 2)^2} = \dfrac{x^2 - 6 - (x^2 - 5x + 6)}{(x - 2)^2} = \mathbf{\dfrac{5x - 12}{(x - 2)^2}}$
            \item $\dfrac{x + 1}{x^2 - 4} \cdot \dfrac{x - 2}{x^2 - 1} = \dfrac{(x + 1)(x - 2)}{(x - 2)(x + 2)(x - 1)(x + 1)} = \mathbf{\dfrac{1}{(x + 2)(x - 1)}}$
            \item $\dfrac{3}{x + 2} + \dfrac{2}{x - 2} - \dfrac{8}{x^2 - 4} = \dfrac{3(x - 2) + 2(x + 2) - 8}{(x - 2)(x + 2)} = \dfrac{3x - 6 + 2x + 4 - 8}{x^2 - 4} = \mathbf{\dfrac{5x - 10}{x^2 - 4}} = \mathbf{\dfrac{5(x - 2)}{(x - 2)(x + 2)}} = \mathbf{\dfrac{5}{x + 2}}$
        \end{enumerate}

    \item \textbf{Expresiones algebraicas en problemas geométricos:}
        \begin{enumerate}[label=\alph*), itemsep=0.8em]
            \item En un triángulo rectángulo, un cateto mide $14\text{ cm}$. Escribe el perímetro y el área del triángulo en función de su hipotenusa $x$.\\
            \noindent{\color{solcolor}\textbf{Solución:}} El otro cateto mide $\sqrt{x^2 - 14^2} = \sqrt{x^2 - 196}$.
            $$\text{Perímetro: } \mathbf{P(x) = 14 + x + \sqrt{x^2 - 196}}$$
            $$\text{Área: } \mathbf{A(x) = \frac{1}{2} \cdot 14 \cdot \sqrt{x^2 - 196} = 7\sqrt{x^2 - 196}}$$

            \item En un triángulo rectángulo, la hipotenusa mide $20\text{ cm}$. Expresa su perímetro y su área en función de la longitud de uno de sus catetos $y$.\\
            \noindent{\color{solcolor}\textbf{Solución:}} El otro cateto es $\sqrt{20^2 - y^2} = \sqrt{400 - y^2}$.
            $$\text{Perímetro: } \mathbf{P(y) = 20 + y + \sqrt{400 - y^2}}$$
            $$\text{Área: } \mathbf{A(y) = \frac{1}{2}y\sqrt{400 - y^2}}$$

            \item Un rectángulo tiene una altura que mide $6\text{ cm}$. Expresa su diagonal y su área en función de su base $x$.\\
            \noindent{\color{solcolor}\textbf{Solución:}} 
            $$\text{Diagonal: } \mathbf{d(x) = \sqrt{x^2 + 6^2} = \sqrt{x^2 + 36}}$$
            $$\text{Área: } \mathbf{A(x) = 6x}$$
        \end{enumerate}

    \item \textbf{Expresa el área de la zona NO edificada en función de las variables indicadas:}
        \begin{enumerate}[label=\alph*), itemsep=1.5em]
            \item En una parcela de lados $x$ e $y$ se construye una casa en la zona indicada...\\
            \noindent{\color{solcolor}\textbf{Solución:}} Área total $= x \cdot y$. Área de la casa $= 30 \cdot (x - 50)$ (según el esquema del gráfico adjunto).
            $$\mathbf{A_{\text{no edificada}}(x, y) = x \cdot y - (y-30)(x - 50)}$$

            \item En un terreno rectangular de ancho $x$ y largo $y$, se reserva una franja exterior no edificable de $4\text{ m}$ de ancho en todo el perímetro. Expresa el área edificable interior en función de $x$ e $y$.\\
            \noindent{\color{solcolor}\textbf{Solución:}} Las dimensiones interiores disminuyen $4\text{ m}$ por cada lado ($8\text{ m}$ en total por dimensión):
            $$\mathbf{A_{\text{edificable}}(x, y) = (x - 8)(y - 8)}$$

            \item En una parcela cuadrada de lado $x$, se edifica una piscina rectangular de dimensiones $12\text{ m}$ y $y\text{ m}$. Expresa el área libre ajardinada en función de $x$ e $y$.\\
            \noindent{\color{solcolor}\textbf{Solución:}} Área total del cuadrado $= x^2$. Área de la piscina $= 12y$.
            $$\mathbf{A_{\text{ajardinada}}(x, y) = x^2 - 12y}$$
        \end{enumerate}

    \item \textbf{Jugando con un trozo de alambre de 36 cm, doblamos el alambre para formar un rectángulo con base "$x$" cm. Hállase la expresión algebraica que define el área del rectángulo en función de $x$ y calcula su valor para $x = 5$.}\\
    \noindent{\color{solcolor}\textbf{Solución:}} 
    El perímetro total es $36\text{ cm}$, por lo que la suma de la base y la altura es la mitad del perímetro: $18\text{ cm}$.
    Si la base es $x$, la altura será $18 - x$.
    $$\text{Área en función de } x\text{: } \mathbf{A(x) = x(18 - x) = 18x - x^2}$$
    Para $x = 5\text{ cm}$:
    $$\mathbf{A(5)} = 5(18 - 5) = 5 \cdot 13 = \mathbf{65\text{ cm}^2}$$

\end{enumerate}
\end{document}